\section{第~\ref{sec:nora}~章：\nt{NORA}}

\nt{NORA}系统的构造、它的两种模式、与瞳孔的联系（不仅经由\nt{NORA}）、信号与背景之比的提高以及行为中的倒U形都是直接测量的；乘法增益$g$是模型；增益切换选择模式并起逆温度的作用，是本章的推论；调节增益的收益是按本书模型的计算；“\nt{NORA}是探索的旋钮”和“\nt{NORA}是中断”是假说。

{\footnotesize
\setlength{\tabcolsep}{3pt}\renewcommand{\arraystretch}{1.15}
\begin{longtable}{>{\raggedright}p{0.5cm}>{\raggedright}p{4.0cm}>{\raggedright}p{5.6cm}>{\raggedright}p{2.3cm}>{\raggedright\arraybackslash}p{1.7cm}}
\toprule
序 & 命题 & 实验与测量内容 & 来源 & 状态 \\
\midrule\endfirsthead
\toprule
序 & 命题 & 实验与测量内容 & 来源 & 状态 \\
\midrule\endhead
1 & 人类的\nt{NORA}神经元有几万个，几乎都在一个核团中 & 人脑干连续切片，酪氨酸羟化酶染色：蓝斑中有$53\,900$个细胞，相邻亚核中另有$6\,260$个。 & \cite{Baker1989LC} & 已测量 \\
2 & 输出几乎散布到全脑，但核团的不同部分在一定程度上服务于不同区域 & 小鼠的病毒—遗传标记：蓝斑的\nt{NORA}神经元接收来自脑内广泛区域的输入，并向脑和脊髓的广泛区域发出输出。大鼠：投射到杏仁核的一组增强恐惧的习得，而投射到前额叶皮层的另一组增强恐惧的消退。 & \cite{Schwarz2015LC, Uematsu2017LC, Poe2020} & 已测量 \\
3 & 紧张性模式：平稳的脉冲，频率从每秒零点几个到几个，随觉醒程度变化 & 大鼠，在睡眠—觉醒周期中记录蓝斑神经元：清醒时频率最高，慢波睡眠中较低，快速眼动睡眠中几乎为零；频率的变化先于状态的转换。 & \cite{AstonJonesBloom1981} & 已测量 \\
4 & 时相性模式：对任务中的重要事件产生短促的簇放电，大约在决策时刻 & 猴子，寻找稀有目标的任务：所有蓝斑神经元只对目标以簇放电响应，在目标出现后$\sim90\ms$、手部反应前$\sim200\ms$；表现差时簇放电较弱。在二选一任务中，簇放电与反应的时间关系比与刺激的更紧密，正确和错误反应时都有。 & \cite{AstonJones1994vigilance, Clayton2004LC} & 已测量 \\
5 & 瞳孔直径是\nt{NORA}的不精确测试点：它既跟随\nt{NORA}，也跟随\nt{ACET}和其他区域 & 猴子：蓝斑的脉冲，直至单个细胞的单个脉冲，都能预测瞳孔扩大；上丘和扣带皮层中也有类似但更晚的联系。小鼠：瞳孔的波动既跟随皮层中去甲肾上腺素能输出的活动，也跟随胆碱能输出的活动。 & \cite{Joshi2016, Reimer2016} & 已测量 \\
6 & 去甲肾上腺素对背景活动的抑制强于对诱发活动的抑制；乘法增益~\eqref{eq:gain}是再现这一效应的模型 & 清醒猴子，听觉皮层，离子电泳去甲肾上腺素：神经元背景活动受到的抑制强于对同类叫声的响应——信噪比升高。乘法S形函数~\eqref{eq:gain}取自一个网络模型；斜率被真的乘上系数，没有测量过。 & \cite{Foote1975NE, ServanSchreiber1990} & 已测量；$g$为模型 \\
7 & 增益只对放大器之后的噪声提高$d'$~\eqref{eq:gainsnr}（命题~\ref{prop:gainnoise}） & 本章推导；在网络模型中增益改善信号检测。没有把放大器前后的噪声分开并检验这一区别的实验。 & 本章推导；\cite{ServanSchreiber1990} & 理论 \\
8 & 边界$g\beta=1$把选择网络从“斟酌”切换到“已决断”~\eqref{eq:wtagain}（命题~\ref{prop:gainwta}，图~\ref{fig:pitchfork}） & 两个相互抑制的组的线性分析；本章推导。脑中没有对这一阈值的直接测量。 & 本章推导 & 理论 \\
9 & 增益是选择的逆温度~\eqref{eq:softmax} & 放大器之后的噪声服从逻辑斯谛分布时，选择概率是$T=\sigma/g$的softmax；本章推导。 & 本章推导 & 理论 \\
10 & \nt{NORA}决定探索多少：高紧张性活动对应小的有效$g$（探索），决策时刻的时相性簇放电对应大的（决断） & 人类：在随机选择之前的基线瞳孔比选择已知最佳项之前的大；瞳孔较大的人探索倾向更高。大鼠：转入随机选择模式受\nt{NORA}向前扣带皮层的输入控制。簇放电提高有效$g$，没有直接测量。 & \cite{JepmaNieuwenhuis2011, Tervo2014stochastic, AstonJones2005} & 假说 \\
11 & 奖励与$g$呈倒U形关系；在快速变化的世界中最佳$g$更小（命题~\ref{prop:explore}，图~\ref{fig:invertedU}） & \texttt{models/nora\_explore.py}，$60$次运行，每次$4000$次尝试。每$1000$次尝试变化一次：最佳$g=16$，比例$0.976$；每$20$次：$g=8$，比例$0.886$；$g=0.5$时为$0.77$。重算结果与本章一致。 & \texttt{models/\allowbreak nora\_\allowbreak explore.py} & 模型 \\
12 & 行为中的倒U形：紧张性活动适中、簇放电清晰时表现最好 & 猴子，与第~4~行相同的任务：表现好的时期紧张性频率适中，对目标的簇放电清晰；紧张性频率高时没有簇放电，错误反应更多。紧张性和诱发活动的变化跟随表现质量的波动。 & \cite{Usher1999LC, AstonJones2005} & 已测量 \\
13 & \nt{NORA}报告非预期的不确定性，\nt{ACET}报告预期的不确定性 & 含两种不确定性的贝叶斯推断理论；所引文献中没有按递质区分这两种信号的直接实验。 & \cite{YuDayan2005} & 假说 \\
14 & 在突然变化之后人学得更快，瞳孔扩大 & 人类，带突变的预测任务：瞳孔的短暂变化跟随对变化概率的估计，并与基线瞳孔一起预测新数据对估计的改变程度（学习速率）；与光照和任务无关的瞳孔变化会改变这一速率。瞳孔在这里显示的正是\nt{NORA}，是根据第~5~行间接数据的推论。 & \cite{Nassar2012} & 已测量 \\
15 & 时相性簇放电起中断作用：网络复位状态 & 没有\nt{NORA}簇放电使网络状态复位的直接实验；只测得了对重要事件的簇放电（第~4~行）。 & \cite{BouretSara2005} & 假说 \\
16 & 按意外调节$g$或学习速率，几乎不比最佳恒定设置更好 & \texttt{models/nora\_explore.py}，每个时期$1000$次尝试的世界（每$1000$次和每$20$次变化一次）。最佳恒定$g=8$：$0.925$；调节$g$：最高$0.927$；调节学习速率：最高$0.926$；恒定速率$0.4$：$0.928$。重算结果与本章一致。 & \texttt{models/\allowbreak nora\_\allowbreak explore.py} & 模型 \\
\bottomrule
\end{longtable}}
